% Slide 2 — Problem Statement
\begin{frame}{Problem Statement}
  \begin{columns}[T,onlytextwidth]
    \column{0.52\textwidth}
    \begin{block}{Motivation \& goal}
      Brain tumor type strongly affects treatment decisions, but manual MRI interpretation
      is time-consuming and subject to variability. We target a reproducible,
      CPU-friendly pipeline suitable for decision support.\vspace{0.3em}

      \textbf{Research question:} Can optimized handcrafted features + classical ML reach strong
      performance on a 3-class brain-tumor MRI task while remaining interpretable?
    \end{block}

    \column{0.46\textwidth}
    \begin{block}{What we do (this talk)}
      \scriptsize
      \begin{tabular}{@{}p{0.30\linewidth}p{0.66\linewidth}@{}}
        \toprule
        \textbf{Step} & \textbf{Outcome} \\
        \midrule
        Feature extraction & Statistics, LBP, DWT, HOG descriptors \\
        Feature optimization & Best params via separability metrics (no training) \\
        Benchmarking & Tuned models across feature sets (\NModels{} $\times$ \NFeatureSets{}) \\
        Validation & CV, bootstrap CI, McNemar / Friedman tests \\
        \bottomrule
      \end{tabular}
    \end{block}

    \vspace{-0.2em}
    {\scriptsize \textit{Scope:} no MRI theory deep dive; no algorithm derivations; focus on methodology and results.}
  \end{columns}
\end{frame}

\note{
Tumor type drives treatment. Manual reading is costly. We focus on an auditable,
CPU-friendly ML pipeline rather than deep learning.
}
